%-------------------------------------------------------------------------------

%%%%%%%%%%%%%%%%%%%%%%%%%%%%%%%%%%%%%%%%%
% Beamer Presentation
% LaTeX Template
% Version 1.0 (10/11/12)
%
% This template has been downloaded from:
% http://www.LaTeXTemplates.com
%
% License:
% CC BY-NC-SA 3.0 (http://creativecommons.org/licenses/by-nc-sa/3.0/)
%
%%%%%%%%%%%%%%%%%%%%%%%%%%%%%%%%%%%%%%%%%

%-------------------------------------------------------------------------------
%	PACKAGES AND THEMES
%-------------------------------------------------------------------------------

\documentclass{beamer}
\usepackage{xcolor}
\usepackage{graphicx}
\usepackage{tikz}
\usepackage{listings}
\usepackage{multicol}

\usepackage{amsmath}

\definecolor{applegreen}{rgb}{0.55, 0.71, 0.0}
\definecolor{blue(ncs)}{rgb}{0.0, 0.45, 0.60}
\definecolor{burgundy}{rgb}{0.5, 0.0, 0.13}

\definecolor{cadet}{rgb}{0.33, 0.41, 0.47}
\definecolor{airforceblue}{rgb}{0.36, 0.54, 0.66}

\lstdefinestyle{C}{
  language=C,
  emptylines=1,
  breaklines=true,
  basicstyle=\ttfamily\color{black},
  identifierstyle=\ttfamily,
  keywordstyle=\color[rgb]{0.0, 0.0, 1.0},
  stringstyle=\color[rgb]{1.0, 0.0, 0.0},
  commentstyle=\color{gray}\slshape,
}

\lstdefinestyle{assembly}{
  emptylines=1,
  breaklines=true,
  basicstyle=\ttfamily\color{black},
  identifierstyle=\ttfamily,
  keywordstyle=\color[rgb]{1.0, 0.0, 1.0},
  stringstyle=\color[rgb]{1.0, 0.0, 0.0},
  commentstyle=\color{gray}\slshape,
  morekeywords={vmovapd, vfmadd231pd,vfnmadd213pd, vmulpd},
}

\lstdefinestyle{qfunc}{
  language=C,
  emptylines=1,
  breaklines=true,
  basicstyle=\ttfamily\color{black},
  commentstyle=\color{gray}\slshape,
  moredelim=**[is][\color{applegreen}]{@}{@},
  moredelim=**[is][\color{blue(ncs)}]{!}{!},
  moredelim=**[is][\color{green!40!black}]{z}{z},
  moredelim=**[is][\color{blue(ncs)!80}]{#}{#},
  moredelim=**[is][\color{blue(ncs)!80!black}]{'}{'},
}

\lstdefinestyle{oper}{
  language=C,
  emptylines=1,
  breaklines=true,
  basicstyle=\ttfamily\color{black},
  moredelim=**[is][\color{applegreen}]{@}{@},
  moredelim=**[is][\color{blue(ncs)}]{!}{!},
  moredelim=**[is][\color{burgundy}]{#}{#},
  moredelim=**[is][\color{red}]{'}{'},
  moredelim=**[is][\bf]{z}{z},
}

\mode<presentation> {

\usetheme{CambridgeUS}

\usecolortheme{wolverine}

\definecolor{gold}{HTML}{D4A017}
\definecolor{darkgold}{HTML}{B7950B}

\setbeamercolor{palette primary}{bg=cadet,fg=white}
\setbeamercolor{palette secondary}{bg=airforceblue,fg=white}
\setbeamercolor{palette tertiary}{bg=black,fg=white}
\setbeamercolor{palette quaternary}{bg=cadet,fg=white}

\setbeamercolor{frametitle}{bg=airforceblue,fg=white}

\setbeamercolor{section number projected}{bg=black,fg=cadet}
\setbeamercolor{item}{fg=black,bg=cadet}

\setbeamertemplate{page number in head/foot}[framenumber]
}

\usepackage{graphicx} % Allows including images
\usepackage{booktabs} % Allows the use of \toprule, \midrule and \bottomrule in tables

%-------------------------------------------------------------------------------
%	TITLE PAGE
%-------------------------------------------------------------------------------

\title[libCEED Finite Element Library]{Motivating Matrix-Free Finite Elements} % The short title appears at the bottom of every slide, the full title is only on the title page

\author{Jeremy L Thompson} % Your name
\institute[CU Boulder] % Your institution as it will appear on the bottom of every slide, may be shorthand to save space
{University of Colorado Boulder \\ % Your institution for the title page
\medskip
\textit{jeremy@jeremylt.org} % Your email address
}
\date{19 Mar 2026} % Date, can be changed to a custom date

\begin{document}

\begin{frame}
\titlepage % Print the title page as the first slide
\end{frame}

%-------------------------------------------------------------------------------

\begin{frame}
\begin{center}
\frametitle{Overview}

Discretization of PDE for computation presents several choices\\

~\\

~\\

Matrix-free finite elements present benefits and challenges\\

~\\

~\\

We'll see how matrix-free finite elements works and the\\
downstream choices and challenges that this discretization provides

\end{center}
\end{frame}
 
%-------------------------------------------------------------------------------

\begin{frame}
\frametitle{Overview} % Table of contents slide, comment this block out to remove it
\tableofcontents
\end{frame}

%-------------------------------------------------------------------------------
%	PRESENTATION SLIDES
%-------------------------------------------------------------------------------

%-------------------------------------------------------------------------------
\section{Finite Elements}
%-------------------------------------------------------------------------------

\begin{frame}
\begin{center}
\frametitle{End Goal}

We need to transform equations:\vspace{-10px}

\begin{equation*}
\begin{split}
- \nabla \cdot \sigma - \rho g &= 0,\, \text{in}\hspace{0.5em} \Omega\\[-5px]
                             u &= \bar{u},\, \text{on}\hspace{0.5em} \partial \Omega^D\\[-5px]
                \sigma \cdot n &= \bar{t},\, \text{on}\hspace{0.5em} \partial \Omega^N
\end{split}
\end{equation*}

\vspace{10px}

into simulations:

\includegraphics[height=4.5cm]{SolidTwistExample}

\end{center}
\end{frame}

%-------------------------------------------------------------------------------

\begin{frame}
\begin{center}
\frametitle{Equation to Solve}

Consider the general non-linear residual

\vspace{-10px}

\begin{equation*}
F \left( u \right) - R = 0
\end{equation*}

\vspace{10px}

where

\begin{itemize}

\item $F \left( u \right)$ represents our 'physics`

\item $R$ is some 'right hand side' or forcing term

\item $u$ is the unknown solution

\end{itemize}

~\\

\end{center}
\end{frame}

%-------------------------------------------------------------------------------

\begin{frame}
\begin{center}
\frametitle{Problem}

Computers like discrete equations/operations\\

~\\

~\\

Our equation is continuous, and computers \emph{hate} that

\end{center}
\end{frame}

%-------------------------------------------------------------------------------

\begin{frame}
\begin{center}
\frametitle{Another Problem}

Solving these non-linear PDEs can be hard\\

~\\

~\\

Really, really hard

\end{center}
\end{frame}

%-------------------------------------------------------------------------------

\begin{frame}
\begin{center}
\frametitle{First Decision}

With the $L^2$ inner product

\begin{equation*}
\left\langle v, u \right\rangle = \int_\Omega v u \, dV
\end{equation*}

We can reformulate to the 'weak form':

\begin{equation*}
\left\langle v, f \left( u \right) \right\rangle = \int_\Omega v \cdot f_0 \left( u, \nabla u \right) + \nabla v : f_1 \left( u, \nabla u \right) \, dV = 0
\end{equation*}

Note:

\begin{itemize}

\item boundary terms are omitted here (integration by parts)

\item $v$ is in some homogeneous space $V_0$

\end{itemize}

\end{center}
\end{frame}

%-------------------------------------------------------------------------------

\begin{frame}
\begin{center}
\frametitle{$f_0$ and $f_1$ Examples}

Projection problem:

\vspace{-10px}

\begin{equation*}
\begin{split}
u &= g\\
\int_\Omega v \cdot u \, dV &= \int_\Omega v g \, dV\\
f_0 \left( u, \nabla u \right) &= u,\, f_1 \left( u, \nabla u \right) = 0
\end{split}
\end{equation*}

~\\

~\\

Poisson problem:

\vspace{-10px}

\begin{equation*}
\begin{split}
-\nabla^2 u &= g\\
\int_\Omega \nabla v : \nabla u \, dV - \int_{\partial \Omega} v \nabla u \cdot n \, dS &= \int_\Omega v g \, dV \\
f_0 \left( u, \nabla u \right) &= 0,\, f_1 \left( u, \nabla u \right) = \nabla u
\end{split}
\end{equation*}


\end{center}
\end{frame}

%-------------------------------------------------------------------------------

\begin{frame}
\begin{center}
\frametitle{Second Decision}

How to discretize?

\begin{itemize}

\item $u, v \in V$\\

~\\

\item Different $V$ give different properties\\

~\\

\item We select $V$ to be the space of piecewise continuous polynomials\\

~\\

\item No continuity in derivatives (not splines)

\end{itemize}

\end{center}
\end{frame}

%-------------------------------------------------------------------------------

\begin{frame}
\begin{center}
\frametitle{Second Decision}

\includegraphics[height=4.7cm]{libCEEDAPI}

T(otal)-vector represents discretization\\

~\\

Want solution values on red dots (Degrees of Freedom [DoF])\\

~\\

Fitting piecewise polynomial to red dots gives continuous solution

\end{center}
\end{frame}

%-------------------------------------------------------------------------------

\begin{frame}
\begin{center}
\frametitle{Third Decision}

Need to convert integration to a discrete problem\\

~\\

\begin{itemize}

\item Quadrature rules take values at points to compute integral\\

~\\

\item Piecewise nature means each element can be integrated separately\\

~\\

\item Workflow:

\begin{enumerate}

\item Map to interpolating polynomial

\item Evaluate values/derivatives at quadrature points

\item Evaluate $f_0$, $f_1$

\item Multiply by quadrature weights

\item Transpose process to return to DoFs

\end{enumerate}

\end{itemize}

\end{center}
\end{frame}

%-------------------------------------------------------------------------------

\begin{frame}
\begin{center}
\frametitle{Third Decision}

\includegraphics[height=4.7cm]{GaussVsTrap}

Why the extra work of interpolating to quadrature points?\\

~\\

Better accuracy in discrete integration\\

~\\

Gauss-Legendre accurate for polynomials up to $2n - 1$

\end{center}
\end{frame}

%-------------------------------------------------------------------------------

\begin{frame}
\begin{center}
\frametitle{Third Decision}

\includegraphics[height=4.7cm]{libCEEDAPI}

E(lement)-vector represents each 'piece' of our solution polynomial\\

~\\

DoFs on blue dots are red dots on a specific element\\

~\\

{\color{blue(ncs)}$B$} maps to the quadrature points - green dots

{\color{applegreen}$D$} evaluates $f_0$, $f_1$, multiply by quadrature weights

\end{center}
\end{frame}

%-------------------------------------------------------------------------------

\begin{frame}
\begin{center}
\frametitle{Galerkin Form}

Galerkin form for an arbitrary second order PDE:\\

\begin{equation*}
\sum_e \mathcal{E}^T \left[ \left( {\color{blue(ncs)}\mathbf{N}}^e \right)^T {\color{applegreen}\mathbf{W}}^e \Lambda \left( {\color{applegreen}f_0} \left( u^e, \nabla u^e \right) \right) + \sum_{i = 0}^{d - 1} \left( {\color{blue(ncs)}\mathbf{D}}_i^e \right)^T {\color{applegreen}\mathbf{W}}^e \Lambda \left( {\color{applegreen}f_1} \left( u^e, \nabla u^e \right) \right) \right] = 0
\label{eq:galerkin_form}
\end{equation*}

\begin{itemize}

\item $\mathcal{E}$ - element assembly/restriction operator\\

\item ${\color{blue(ncs)}\mathbf{N}}^e$ - interpolation to quadrature points\\

\item ${\color{blue(ncs)}\mathbf{D}}_i^e$ - derivatives at quadrature points\\

\item ${\color{applegreen}\mathbf{W}}^e$ - quadrature weights\\

\item $\Lambda$ - pointwise multiplication at quadrature points\\

\item $u^e = {\color{blue(ncs)}\mathbf{N}}^e \mathcal{E}^e u$ and $\nabla u^e = \lbrace {\color{blue(ncs)}\mathbf{D}}_i^e \mathcal{E}^e u \rbrace_{i = 0}^{d - 1}$

\end{itemize}

\end{center}
\end{frame}

%-------------------------------------------------------------------------------

\begin{frame}
\begin{center}
\frametitle{Wait, What?!?}

Non-linear problem:\\

\vspace{-20px}

\begin{equation*}
F \left( u \right) - R = 0
\end{equation*}

~\\

Weak form:\\

\vspace{-20px}

\begin{equation*}
\left\langle v, f \left( u \right) \right\rangle = \int_\Omega v \cdot f_0 \left( u, \nabla u \right) + \nabla v : f_1 \left( u, \nabla u \right) \, dV = 0
\end{equation*}

~\\

Galerkin form:\\

\vspace{-20px}

\begin{equation*}
\sum_e \mathcal{E}^T \left[ \left( {\color{blue(ncs)}\mathbf{N}}^e \right)^T {\color{applegreen}\mathbf{W}}^e \Lambda \left( {\color{applegreen}f_0} \left( u^e, \nabla u^e \right) \right) + \sum_{i = 0}^{d - 1} \left( {\color{blue(ncs)}\mathbf{D}}_i^e \right)^T {\color{applegreen}\mathbf{W}}^e \Lambda \left( {\color{applegreen}f_1} \left( u^e, \nabla u^e \right) \right) \right] = 0
\end{equation*}

\end{center}
\end{frame}

%-------------------------------------------------------------------------------

\begin{frame}
\begin{center}
\frametitle{What I'm Handwaving}

\begin{itemize}

\item Elements can have deformed shapes (unstructured meshes)\\

~\\

\item Need 'geometric' factors for each individual element\\
(Change of coordinate systems from physical to reference)\\

~\\

\item ${\color{burgundy}\mathcal{P}}$ is parallel communication (supercomputers!)\\

~\\

\item $R$ (and its weak analogue) can be considered part of $f_0$ or separate\\

~\\

\item Boundary terms are similar but 1 dimension lower (i.e. 3D $\rightarrow$ 2D)

\end{itemize}

\end{center}
\end{frame}

%-------------------------------------------------------------------------------
\section{Matrix-Free}
%-------------------------------------------------------------------------------

\begin{frame}
\begin{center}
\frametitle{Traditional Approach}

\begin{itemize}

\item The Galerkin form naturally creates a sparse matrix for linear $F \left( u \right)$\\

~\\

\item This sparse matrix can be solved with linear algebra\\

~\\

\item Two problems:\\

~\\

\begin{enumerate}

\item $F \left( u \right)$ is not usually linear\\

~\\

\item Sparse matrices don't use hardware optimally\\

\end{enumerate}

~\\

\item Big Idea: perform the operations without forming a matrix

\end{itemize}

\end{center}
\end{frame}

%-------------------------------------------------------------------------------

\begin{frame}
\begin{center}
\frametitle{Modern Hardware}

\includegraphics[height=5.5cm]{peakFlopsAndBandwidth_tall}
\hspace{1cm}
\includegraphics[height=5.5cm]{peakRatio_tall}

Modern hardware has lower memory bandwidth than FLOPs

\end{center}
\end{frame}

%-------------------------------------------------------------------------------

\begin{frame}
\begin{center}
\frametitle{Matrix-Free}

\includegraphics[height=5.5cm]{assembledVsMatrixFree_tall}
\hspace{1cm}
\includegraphics[height=5.5cm]{assembledVsMatrixFreeBalance_tall}

{\small Requirements for matrix-vector product with sparse matrix vs matrix-free\\ for screened Poisson $\nabla^2 u - \alpha^2 u = f$ in 3D}\\

{\bf Matrix-free representations using tensor product bases\\better match modern hardware limitations}

\end{center}
\end{frame}

%-------------------------------------------------------------------------------
\section{Non-Linear Problem}
%-------------------------------------------------------------------------------

\begin{frame}
\begin{center}
\frametitle{Non-linear Problem}

We still have a non-linear problem\\

~\\

~\\

Linear algebra solves \emph{linear} problems\\

\end{center}
\end{frame}

%-------------------------------------------------------------------------------

\begin{frame}
\begin{center}
\frametitle{Newton's Method}

Many non-linear solvers fit into the 'Newton' family\\

\includegraphics[height=4.5cm]{NewtonIteration}

In Newton's method, we form a tangent line\\

~\\

The \emph{line}ar equation is easy to solve\\

~\\

In higher dimensions, the line becomes a (hyper)plane

\end{center}
\end{frame}

%-------------------------------------------------------------------------------

\begin{frame}
\begin{center}
\frametitle{Newton's Method}

We want to solve

\begin{equation*}
F \left( u \right) - R = 0
\end{equation*}

\vspace{10px}

We can instead solve a series of 'tangent plane' problems

\vspace{-10px}

\begin{equation*}
F \left( u_k \right) - R + J \left( u_k \right) \left( u_k - u_{k + 1} \right) = 0
\end{equation*}

$u_{k + 1}$ provides closer estimates of the true solution $u_*$\\

~\\

We'll denote $w = u_k - u_{k + 1}$

\end{center}
\end{frame}

%-------------------------------------------------------------------------------

\begin{frame}
\begin{center}
\frametitle{A New Problem}

Ok, but we still need to form the tangent plane $J \left( u_k \right)$\\

~\\

~\\

This feels like we keep changing what the problem is!

\end{center}
\end{frame}

%-------------------------------------------------------------------------------

\begin{frame}
\begin{center}
\frametitle{Beauty}

Weak form:\\

\begin{equation*}
\left\langle v, f \left( u \right) \right\rangle = \int_\Omega v \cdot f_0 \left( u, \nabla u \right) + \nabla v : f_1 \left( u, \nabla u \right) \, dV = 0
\end{equation*}

We want to take the derivative of this weak form\\

~\\

Note:

\begin{itemize}

\item $\frac{d}{dx} \left( f \left( x \right) + g \left( x \right) \right) = f' \left( x \right) + g' \left( x \right)$\\

~\\

\item $\frac{d}{dx} f \left( g \left( x \right) \right) = f' \left( g \left( x \right) \right) g' \left( x \right)$

\end{itemize}

\end{center}
\end{frame}

%-------------------------------------------------------------------------------

\begin{frame}
\begin{center}
\frametitle{Beauty}

Weak form derivative:\\

\begin{equation*}
\langle v, J \left( u_k \right) w \rangle = \int_{\Omega}
\left[ \begin{array}{c c}
v^T & \left( \nabla v \right)^T
\end{array} \right]
\left[ \begin{array}{c c}
f_{0, 0} & f_{0, 1}\\
f_{1, 0} & f_{1, 1}
\end{array} \right]
\left[ \begin{array}{c}
w \\ \nabla w
\end{array} \right]
\label{eq:jacobian_form}
\end{equation*}

~\\

Note:

\begin{itemize}

\item $f_{i, 0} = \frac{\partial f_i}{\partial u} \left( u_k, \nabla u_k \right)$\\

~\\

\item $f_{i, 1} = \frac{\partial f_i}{\partial \nabla u} \left( u_k, \nabla u_k \right)$\\

~\\

\item Can be discretized the same way as before

\end{itemize}

\end{center}
\end{frame}

%-------------------------------------------------------------------------------

\begin{frame}
\begin{center}
\frametitle{Wait, What?!? - Non-linear Equation}

Non-linear problem:\\

\vspace{-20px}

\begin{equation*}
F \left( u \right) - R = 0
\end{equation*}

~\\

Weak form:\\

\vspace{-20px}

\begin{equation*}
\left\langle v, f \left( u \right) \right\rangle = \int_\Omega v \cdot f_0 \left( u, \nabla u \right) + \nabla v : f_1 \left( u, \nabla u \right) \, dV = 0
\end{equation*}

~\\

Galerkin form:\\

\vspace{-20px}

\begin{equation*}
\sum_e \mathcal{E}^T \left[ \left( {\color{blue(ncs)}\mathbf{N}}^e \right)^T {\color{applegreen}\mathbf{W}}^e \Lambda \left( {\color{applegreen}f_0} \left( u^e, \nabla u^e \right) \right) + \sum_{i = 0}^{d - 1} \left( {\color{blue(ncs)}\mathbf{D}}_i^e \right)^T {\color{applegreen}\mathbf{W}}^e \Lambda \left( {\color{applegreen}f_1} \left( u^e, \nabla u^e \right) \right) \right] = 0
\end{equation*}

\end{center}
\end{frame}

%-------------------------------------------------------------------------------

\begin{frame}
\begin{center}
\frametitle{Wait, What?!? - Linear Equation}

Linearized problem:\\

\vspace{-20px}

\begin{equation*}
F \left( u_k \right) - R + J \left( u_k \right) \left( u_k - u_{k + 1} \right) = 0
\end{equation*}

~\\

Weak form:\\

\vspace{-20px}

\begin{equation*}
\langle v, J \left( u_k \right) w \rangle = \int_{\Omega}
\left[ \begin{array}{c c}
v^T & \left( \nabla v \right)^T
\end{array} \right]
\left[ \begin{array}{c c}
f_{0, 0} & f_{0, 1}\\
f_{1, 0} & f_{1, 1}
\end{array} \right]
\left[ \begin{array}{c}
w \\ \nabla w
\end{array} \right]
\end{equation*}

~\\

Galerkin form:\\

\vspace{-20px}

\begin{equation*}
\sum_e \mathcal{E}^T \left[ \left( {\color{blue(ncs)}\mathbf{N}}^e \right)^T {\color{applegreen}\mathbf{W}}^e \Lambda \left( {\color{applegreen}f^*_0} \left( w^e, \nabla w^e \right) \right) + \sum_{i = 0}^{d - 1} \left( {\color{blue(ncs)}\mathbf{D}}_i^e \right)^T {\color{applegreen}\mathbf{W}}^e \Lambda \left( {\color{applegreen}f^*_1} \left( w^e, \nabla w^e \right) \right) \right] = 0
\end{equation*}

\end{center}
\end{frame}

%-------------------------------------------------------------------------------
\section{Preconditioning}
%-------------------------------------------------------------------------------

\begin{frame}
\begin{center}
\frametitle{A New Problem}

We have a discrete system we can solve efficiently\\

~\\

~\\

But we still have a problem\\

~\\

~\\

(Our problems are getting easier to tackle though!)

\end{center}
\end{frame}

%-------------------------------------------------------------------------------

\begin{frame}
\begin{center}
\frametitle{Condition Number}

Matrix-free representation of the linear problem restricts solvers\\

~\\

~\\

Krylov subspace methods only need the matrix-vector product $A x$,\\
not the assembled matrix $A$\\

~\\

~\\

But Krylov subspace methods are sensitive to condition number\\

~\\

~\\

High-order finite elements are more efficient matrix-free,\\
but high-order finite elements have a higher condition number

\end{center}
\end{frame}

%-------------------------------------------------------------------------------

\begin{frame}
\begin{center}
\frametitle{Preconditioning}

Preconditioning improves the condition number of $A$\\

~\\

~\\

Big Idea: $P^{-1} A \approx I$ has a better condition number

~\\

~\\

Challenge: Finding efficient $P^{-1}$ can difficult\\

~\\

~\\

The Fix: This is a rich research area with lots of options

\end{center}
\end{frame}

%-------------------------------------------------------------------------------
\section{Cool Results}
%-------------------------------------------------------------------------------

\begin{frame}
\begin{center}
\frametitle{Example - Brittle Fracture}

\includegraphics[height=5cm]{compressive-shear-damage-FEM}

~\\

Quasistatic simulation of a compressive shear test for a generic brittle material

\end{center}
\end{frame}

%-------------------------------------------------------------------------------

\begin{frame}
\begin{center}
\frametitle{Example - Schwarz Pendulum}

\includegraphics[height=3.5cm]{dynamic-pendulum}

~\\

Dynamic simulation of 8 Schwarz-P cells with an applied traction force creating pendulum-like motion

\end{center}
\end{frame}

%-------------------------------------------------------------------------------

\begin{frame}
\begin{center}
\frametitle{Example - Sinker Simulation}

\includegraphics[height=3.5cm]{sinker-MPM-FEM}

~\\

Quasistatic simulation of a cylinder with a dense high-modulus inclusion surrounded by a low-density low-modulus near-incompressible "foam"

\end{center}
\end{frame}

%-------------------------------------------------------------------------------

\begin{frame}
\begin{center}
\frametitle{Example - Press Simulation}

\includegraphics[height=5cm]{iMPMCloseup}
\hspace*{1cm}
\includegraphics[height=5cm]{iMPMCompress}

~\\

Compression of mock HE grains (gold) and binder (pink) mixture\\

(Reset background mesh to computational region on each timestep)

\end{center}
\end{frame}

%-------------------------------------------------------------------------------
\section{Questions}
%-------------------------------------------------------------------------------

\begin{frame}
\begin{center}
\frametitle{Questions?}

\begin{center}
\includegraphics[height=2.75cm]{libCEEDlogo.png}
\includegraphics[height=2.75cm]{Ratellogo.png}
\end{center}

{\flushleft

libCEED Repo: \href{https://github.com/CEED/libCEED}{https://github.com/CEED/libCEED}\\
Ratel Repo: \hspace*{0.39cm} \href{https://gitlab.com/micromorph/ratel}{https://gitlab.com/micromorph/ratel}\\

~\\

Grant: Predictive Science Academic Alliance Program (DE-NA0003962)\\

}

\begin{center}

\includegraphics[height=0.8cm]{psaap-center-logos.png}
\end{center}

\end{center}
\end{frame}

%-------------------------------------------------------------------------------

\begin{frame}
\titlepage % Print the title page
\end{frame}

%-------------------------------------------------------------------------------
\end{document}
%-------------------------------------------------------------------------------
